\section{Tecnologias do Back-end}
\label{sec:tecnologias-backend}

\subsection{Flask e Blueprints}
\label{subsec:flask-blueprints}

O Flask foi utilizado na API do módulo por permitir incorporar suas rotas à aplicação existente
sem impor uma reorganização do sistema. Os \textit{blueprints} agruparam as operações de
criação, análise, acompanhamento e reversão das solicitações em um componente isolado,
preservando as demais funcionalidades e facilitando a manutenção do módulo \cite{flask2024}.

\subsection{PostgreSQL}
\label{subsec:postgresql-transacoes}

O módulo utiliza PostgreSQL para armazenar solicitações, alterações cadastrais e registros de
auditoria. A persistência dessas informações precisa manter o cadastro e o histórico coerentes,
sobretudo quando uma decisão altera mais de uma entidade.

O PostgreSQL foi adequado por oferecer transações para as operações compostas do fluxo. Na
aprovação, a mudança de categoria, a atualização da pontuação, a gravação do \textit{snapshot}
do estado anterior e a atualização do status da solicitação devem ser confirmadas em conjunto;
se uma etapa falhar, nenhuma alteração deve permanecer aplicada \cite{postgresql2024}.

\subsection{JSONB e Modelagem de Estado}
\label{subsec:jsonb}

O PostgreSQL oferece o tipo de dado JSONB, que armazena documentos JSON de forma binária,
permitindo consultas e indexação eficientes \cite{postgresql2024}. O módulo utiliza esse recurso
para registrar o estado do especialista no momento da solicitação, preservando a categoria, a
pontuação e outros atributos relevantes. Essa abordagem facilita a reversão de alterações e a
auditoria das mudanças, mantendo o histórico completo das operações realizadas.

% TO-DO: Exemplo visual de um registro JSONB, mostrando a categoria e pontuação antes da 
% alteração, e o status da solicitação.

\todo[inline]{TODO:Exemplo visual de um registro JSONB, mostrando a categoria e pontuação antes da 
alteração, e o status da solicitação.}

\subsection{SQLAlchemy e Unidade de Trabalho}
\label{subsec:sqlalchemy-unidade-trabalho}

O SQLAlchemy foi adotado como ORM para coordenar as alterações de entidades relacionadas em uma
única unidade de trabalho \cite{bayer2012}. A sessão acumula as mudanças e, no \textit{commit},
efetiva-as na transação; em caso de falha, o \textit{rollback} descarta o conjunto. Assim, a
aprovação não deixa a categoria atualizada sem o \textit{snapshot} correspondente ou sem o novo
status da solicitação.

% TO-DO: Diagrama de sequência mostrando a aprovação de uma solicitação, com a atualização da
% categoria, a gravação do snapshot e a alteração do status da solicitação, tudo dentro
% de uma única transação.
 
\todo[inline]{TODO:Diagrama de sequência mostrando a aprovação de uma solicitação, com a atualização da
categoria, a gravação do snapshot e a alteração do status da solicitação, tudo dentro
de uma única transação.}

\subsubsection{Mapeamento Objeto-Relacional e Estratégias de Herança}
\label{subsubsec:mapeamento-objeto-relacional-estrategias-heranca}

O SQLAlchemy oferece diferentes estratégias de herança para mapear hierarquias de classes. A
estratégia Joined Table Inheritance (herança de tabela unida) cria tabelas separadas para a
classe base e suas subclasses, sendo útil quando há muitos atributos exclusivos. Já a estratégia
Single Table Inheritance (herança de tabela única) armazena os diferentes tipos de entidades em
uma única tabela, utilizando uma coluna discriminadora. No módulo de solicitações, adotou-se o
Single Table Inheritance para os tipos de solicitação (mudança de categoria, recurso e
pontuação), pois eles compartilham a maior parte dos atributos, o que evita junções complexas e
otimiza as consultas de listagem.

\subsection{Validação de Dados com Marshmallow e Flasgger}
\label{subsec:validacao-dados-marshmallow-flasgger}

Na fronteira da API, os \textit{schemas} do Marshmallow validam os campos e formatos recebidos
antes que os dados cheguem às regras de negócio \cite{marshmallow2024}. Ao declarar os formatos
esperados para as solicitações e suas respostas, também estabelecem o contrato consumido pelas
interfaces de criação, acompanhamento e análise do módulo.
